\documentclass[
  spanish,
  aspectratio=169,
  xcolor={dvipsnames,table}
]{beamer}

\geometry{paperwidth=19.2cm,paperheight=10.8cm}

\usepackage{fontspec}
\setmainfont{Arial}
\setsansfont{Arial}
\usepackage[spanish,es-tabla,es-noshorthands]{babel}
\usepackage{array}
\usepackage{booktabs}
\usepackage{graphicx}
\usepackage{ragged2e}
\usepackage{tabularx}
\usepackage{amsmath,amssymb}
\usepackage{tikz}
\usepackage{hyperref}

\definecolor{uanlblue}{RGB}{0,55,108}
\definecolor{uanlgold}{RGB}{218,165,32}
\definecolor{tableblue}{RGB}{233,240,247}
\definecolor{darkcharcoal}{RGB}{33,37,41}

\setbeamercolor{palette primary}{bg=uanlblue,fg=white}
\setbeamercolor{palette secondary}{bg=uanlgold,fg=uanlblue}
\setbeamercolor{structure}{fg=uanlblue}
\setbeamercolor{title}{fg=white,bg=uanlblue}
\setbeamercolor{frametitle}{fg=uanlblue,bg=white}
\setbeamerfont{frametitle}{series=\bfseries,size=\large}

\setbeamertemplate{navigation symbols}{}
\setbeamertemplate{footline}{
  \leavevmode%
  \hbox{%
  \begin{beamercolorbox}[wd=.333333\paperwidth,ht=2.5ex,dp=1.1ex,left,leftskip=1cm]{palette primary}%
    \fontsize{7}{9}\selectfont UANL \textbar\ Facultad de Agronomía
  \end{beamercolorbox}%
  \begin{beamercolorbox}[wd=.333333\paperwidth,ht=2.5ex,dp=1.1ex,center]{palette primary}%
    \fontsize{7}{9}\selectfont Cálculo Integral — PIA
  \end{beamercolorbox}%
  \begin{beamercolorbox}[wd=.333333\paperwidth,ht=2.5ex,dp=1.1ex,right,rightskip=1cm]{palette primary}%
    \fontsize{7}{9}\selectfont \insertframenumber{} / \inserttotalframenumber
  \end{beamercolorbox}}%
  \vskip0pt%
}

\begin{document}

\begin{frame}[plain]
  \begin{tikzpicture}[remember picture,overlay]
    \fill[uanlblue] (current page.south west) rectangle (current page.north east);
    \fill[uanlgold] ([yshift=1.2cm]current page.south west) rectangle ([yshift=1.45cm]current page.south east);
  \end{tikzpicture}
  \vspace{-0.5cm}
  \begin{center}
    \includegraphics[height=1.8cm]{UANL/assets-uanl/logos-institucionales.pdf}
    \vspace{0.4cm}
    
    {\fontsize{10}{12}\selectfont \textbf{\color{uanlgold}UNIVERSIDAD AUTÓNOMA DE NUEVO LEÓN \textbar\ FACULTAD DE AGRONOMÍA}}\\[0.2cm]
    {\fontsize{15}{18}\selectfont \textbf{\color{white}MODELACIÓN ANALÍTICA DE ALMACENAMIENTO HIDRÁULICO Y BOMBEO MEDIANTE CÁLCULO INTEGRAL}}\\[0.2cm]
    {\fontsize{11}{13}\selectfont \color{white}Producto Integrador de Aprendizaje (PIA) \textbar\ 30 Puntos}\\[0.5cm]

    \begin{tabular}{rl}
      \textbf{\color{uanlgold}Estudiante:} & \color{white}Carlos Mauricio Carvajal Coronado (Matrícula: 1977875) \\
      \textbf{\color{uanlgold}Cátedra:} & \color{white}Cálculo Integral \textbar\ REF\#261013755131 \\
      \textbf{\color{uanlgold}Programa:} & \color{white}Ingeniero Agrónomo \textbar\ Modelo por Competencias
    \end{tabular}
  \end{center}
\end{frame}

\begin{frame}{1. Justificación Hidroagrícola en Nuevo León}
  \begin{columns}[T]
    \begin{column}{0.48\textwidth}
      \textbf{\color{uanlblue}Contexto y Retos Regionales:}
      \begin{itemize}
        \item Escasez hídrica recurrente en el noreste mexicano.
        \item Necesidad de almacenamiento eficiente en bordos de captación parcelaria.
        \item Optimización energética en bombeo presurizado para fertirriego.
      \end{itemize}
    \end{column}
    \begin{column}{0.48\textwidth}
      \textbf{\color{uanlblue}Aplicación del Cálculo Integral:}
      \begin{itemize}
        \item \textbf{Sólidos de revolución:} Geometría exacta del reservorio ($y = 0.05 x^2$, $H = 5\text{ m}$).
        \item \textbf{Trabajo hidrostático:} Integración de masa y gradiente de elevación $W = \int \rho g A(y)(H-y)\,dy$.
        \item \textbf{Integrales impropias:} Persistencia y decaimiento acumulado de nitratos en suelo.
      \end{itemize}
    \end{column}
  \end{columns}
\end{frame}

\begin{frame}{2. Modelado Volumétrico del Reservorio ($V = 250\pi\text{ m}^3$)}
  \begin{columns}[T]
    \begin{column}{0.48\textwidth}
      \textbf{\color{uanlblue}Método de Discos:}
      \begin{itemize}
        \item Eje de giro vertical $y$, radio $x = \sqrt{20y}$.
        \item Elemento diferencial: $dV = \pi (20y)\,dy$.
        \item Integral:
          $$V = \int_{0}^{5} 20\pi y\,dy = 20\pi \left[ \frac{y^2}{2} \right]_{0}^{5} = 250\pi\text{ m}^3$$
        \item Capacidad: $\approx 785.4\text{ m}^3$ ($785,400\text{ litros}$).
      \end{itemize}
    \end{column}
    \begin{column}{0.48\textwidth}
      \textbf{\color{uanlblue}Método de Cascarones Cilíndricos:}
      \begin{itemize}
        \item Radio $x$, altura $h(x) = 5 - 0.05 x^2$.
        \item Elemento diferencial: $dV = 2\pi x (5 - 0.05 x^2)\,dx$.
        \item Integral:
          $$V = 2\pi \int_{0}^{10} (5x - 0.05 x^3)\,dx = 250\pi\text{ m}^3$$
        \item \textbf{Concordancia exacta} entre métodos analíticos independientes.
      \end{itemize}
    \end{column}
  \end{columns}
\end{frame}

\begin{frame}{3. Trabajo Mecánico de Vaciado y Potencia de Bombeo}
  \begin{block}{Formulación del Trabajo Hidrostático}
    $$W = \rho g \int_{0}^{5} \pi (20y)(5 - y)\,dy = 20\pi (9810) \left[ \frac{5y^2}{2} - \frac{y^3}{3} \right]_{0}^{5} = 1.284 \times 10^7\text{ J} \quad (12.84\text{ MJ})$$
  \end{block}
  \begin{itemize}
    \item \textbf{Tiempo de operación requerido:} $T = 4\text{ horas} = 14,400\text{ segundos}$.
    \item \textbf{Potencia útil teórica:} $P_{\text{útil}} = \frac{12.84\text{ MJ}}{14,400\text{ s}} = 891.7\text{ Watts}$.
    \item \textbf{Potencia al freno con $\eta = 75\%$:} $P_{\text{bomba}} = 1,188.9\text{ W} \approx 1.59\text{ HP} \implies$ Motobomba de $2.0\text{ HP}$.
    \item \textbf{Arreglo solar fotovoltaico:} 6 paneles de $450\text{ W}$ en conexión serie-paralelo.
  \end{itemize}
\end{frame}

\begin{frame}{4. Integrales Impropias y Conclusiones del Proyecto}
  \begin{columns}[T]
    \begin{column}{0.48\textwidth}
      \textbf{\color{uanlblue}Persistencia de Nitratos en Suelo:}
      \begin{itemize}
        \item Tasa de degradación: $C(t) = 45 e^{-0.08 t}$.
        \item Exposición integral acumulada:
          $$\int_{0}^{\infty} 45 e^{-0.08 t}\,dt = \frac{45}{0.08} = 562.5\text{ mg}\cdot\text{día/L}$$
        \item \textbf{Convergencia estricta:} Evidencia ausencia de saturación ecotoxicológica en el manto freático.
      \end{itemize}
    \end{column}
    \begin{column}{0.48\textwidth}
      \textbf{\color{uanlblue}Conclusiones Clave:}
      \begin{enumerate}
        \item Se dimensionó con rigor matemático un reservorio de riego tecnificado.
        \item Se vinculó el cálculo diferencial e integral a la ingeniería de riego solar.
        \item Se cumplieron con excelencia los criterios de evaluación del PIA de Cálculo Integral de la UANL.
      \end{enumerate}
    \end{column}
  \end{columns}
\end{frame}

\end{document}

